\chapter{آرایه‌ها: نگهداریِ چند مقدار}%
\label{ch:arrays}

تا اینجا هر متغیر فقط \emph{یک} مقدار نگه می‌داشت. اما مسئله‌های واقعی با
دسته‌های داده سروکار دارند: نمره‌های یک کلاس، دمای روزهای هفته، امتیازهای
یک مسابقه. با متغیرهای تکی کارمان پیش نمی‌رود؛ برای بیست نمره که نمی‌شود
بیست متغیرِ \code{نمره۱} تا \code{نمره۲۰} ساخت! ابزارِ درست
\textbf{آرایه} است.

\section{آرایه چیست؟}

آرایه ردیفی از جعبه‌های \emph{هم‌گونه} است که زیرِ یک نام کنارِ هم چیده و
شماره‌گذاری شده‌اند:

\begin{salamfa}
روال ریشه:
    نمرات : صحیح[5] = [12, 18, 9, 20, 15]
    سرچاپ نمرات[0]
    سرچاپ نمرات[3]
پایان
\end{salamfa}

\begin{progout}
12
20
\end{progout}

خطِ اول را بشکافیم: \code{صحیح[5]} یعنی «آرایه‌ای از ۵ عددِ صحیح» و
مقدارهای اولیه میانِ کروشه، با کاما جدا شده‌اند.

\subsection{اندیس از صفر شروع می‌شود}

شمارهٔ هر خانه را \textbf{اندیس} می‌گویند و شماره‌گذاری از \emph{صفر} آغاز
می‌شود: \code{نمرات[0]} خانهٔ اول است (۱۲)، \code{نمرات[1]} خانهٔ دوم
(۱۸) و \code{نمرات[4]} خانهٔ آخر (۱۵). حالا رازِ فصلِ \ref{ch:loops} هم
گشوده شد: \code{تکرار 5 از i} شمارنده را از ۰ تا ۴ می‌گرداند چون این‌ها
دقیقاً اندیس‌های یک آرایهٔ ۵تایی‌اند؛ حلقه و آرایه برای هم ساخته شده‌اند.

\begin{warn}
در آرایهٔ ۵تایی، اندیسِ معتبر ۰ تا ۴ است. \code{نمرات[5]} از مرزِ آرایه
بیرون می‌زند و خطاست؛ این از رایج‌ترین اشتباه‌های تمامِ برنامه‌نویسانِ
دنیاست و حتی نامِ مشهوری هم دارد: \emph{خطای یکی کم و زیاد}. هر جا حلقه و
آرایه با هم دیدید، یک لحظه بایستید و بپرسید: اولین و آخرین اندیسی که این
حلقه می‌سازد چیست؟
\end{warn}

خانه‌های آرایه مثلِ متغیر خواندنی و نوشتنی‌اند:

\begin{salamfa}
روال ریشه:
    ناپایا دماها : صحیح[3] = [21, 24, 19]
    دماها[2] = 23
    سرچاپ دماها[0], دماها[1], دماها[2]
پایان
\end{salamfa}

\begin{progout}
21 24 23
\end{progout}

\section{پیمایش: گشتن در آرایه}

مهم‌ترین کاری که با آرایه می‌کنیم \textbf{پیمایش} است: دیدنِ خانه‌ها یکی
پس از دیگری. راهِ همه‌کاره، حلقه با اندیس است؛ و برای اینکه طولِ آرایه را
با عددِ دستی ننویسیم، از \code{len} کمک می‌گیریم که تعدادِ خانه‌ها را
می‌دهد:

\begin{salamfa}
روال ریشه:
    نمرات : صحیح[5] = [12, 18, 9, 20, 15]
    ناپایا مجموع : صحیح = 0
    ناپایا i : صحیح = 0
    تا i < len(نمرات):
        مجموع = مجموع + نمرات[i]
        i = i + 1
    پایان
    سرچاپ "مجموع:", مجموع
    سرچاپ "میانگین:", مجموع / len(نمرات)
پایان
\end{salamfa}

\begin{progoutfa}
مجموع: 74
میانگین: 14
\end{progoutfa}

به شرطِ \code{i < len(نمرات)} دقت کنید: کوچک‌ترِ اکید، نه \code{<=}. با
طولِ ۵، حلقه اندیس‌های ۰ تا ۴ را می‌سازد؛ دقیقاً همهٔ خانه‌ها، نه یکی
بیشتر. این زوجِ «شروع از ۰» و «شرطِ \code{< len}» را مثلِ یک قالبِ حاضری
به خاطر بسپارید.

اگر به اندیس نیاز ندارید و فقط مقدارها را می‌خواهید، سلام راهِ کوتاه‌تر و
خواناتری هم دارد؛ حلقهٔ \code{هر}:

\begin{salamfa}
روال ریشه:
    نمرات : صحیح[5] = [12, 18, 9, 20, 15]
    ناپایا مجموع : صحیح = 0
    هر نمره از نمرات:
        مجموع = مجموع + نمره
    پایان
    سرچاپ "مجموع:", مجموع
پایان
\end{salamfa}

خطِ \code{هر نمره از نمرات:} چنین می‌خواند: «به ازای \emph{هر} خانه از
آرایهٔ نمرات، مقدارش را در \code{نمره} بگذار و بدنه را اجرا کن.» نه
شمارنده‌ای، نه شرطی، نه خطری برای بیرون زدن از مرز.

\section{الگوهای کلاسیکِ آرایه}

سه مسئلهٔ کوچک که ستونِ فقراتِ فصلِ آینده‌اند، همگی با پیمایش و همان الگوی
انباشتن حل می‌شوند.

\subsection{شمردن با شرط}

چند نفر نمرهٔ ۱۷ به بالا گرفته‌اند؟

\begin{salamfa}
روال ریشه:
    نمرات : صحیح[6] = [12, 18, 9, 20, 15, 17]
    ناپایا تعداد : صحیح = 0
    هر نمره از نمرات:
        اگر نمره >= 17:
            تعداد = تعداد + 1
        پایان
    پایان
    سرچاپ "ممتازها:", تعداد
پایان
\end{salamfa}

\begin{progoutfa}
ممتازها: 3
\end{progoutfa}

\subsection{پیدا کردنِ بیشینه}

الگوریتمِ «بزرگ‌ترینِ فعلی» از فصلِ \ref{ch:algorithms} حالا برای هر
اندازه آرایه‌ای کار می‌کند و می‌بینید که ترجمهٔ شبه‌کدِ آن روز چقدر
سرراست است:

\begin{salamfa}
روال بیشینه آرایه(اعداد : صحیح[:]) : صحیح:
    ناپایا بزرگترین : صحیح = اعداد[0]
    هر عدد از اعداد:
        اگر عدد > بزرگترین:
            بزرگترین = عدد
        پایان
    پایان
    برگشت بزرگترین
پایان

روال ریشه:
    نمرات : صحیح[5] = [12, 18, 9, 20, 15]
    سرچاپ "بالاترین نمره:", بیشینه آرایه(نمرات)
پایان
\end{salamfa}

\begin{progoutfa}
بالاترین نمره: 20
\end{progoutfa}

نکتهٔ تازه در سطرِ تعریف است: گونهِ پارامتر \code{صحیح[:]} نوشته شده، یعنی
«دنباله‌ای از اعدادِ صحیح با هر طولی». به این ترتیب همین یک روال برای
آرایهٔ ۵تایی و ۵۰۰تایی یکسان کار می‌کند. دقت کنید که «بزرگ‌ترینِ فعلی» را
با \code{اعداد[0]} شروع کردیم نه با صفر؛ اگر با صفر شروع می‌کردیم، برای
آرایه‌ای که همهٔ اعضایش منفی‌اند جوابِ غلط می‌گرفتیم. حالت‌های مرزی،
همیشه و همه‌جا!

\section{تمرین‌ها}

\exercise آرایه‌ای از دمای ۷ روزِ هفته بسازید. میانگینِ دما و تعدادِ
روزهای گرم‌تر از میانگین را چاپ کنید (دو پیمایشِ جداگانه لازم دارید؛
چرا نمی‌شود در یک پیمایش هر دو را حساب کرد؟).

\exercise روالِ \code{کمینه آرایه(اعداد صحیح[:]) صحیح} را به سبکِ
\code{بیشینه آرایه} بنویسید و بیازمایید.

\exercise آرایه‌ای ۱۰تایی بسازید و با حلقه، خانهٔ \code{i}ام را برابرِ
\code{i * i} کنید؛ سپس همه را چاپ کنید (مربعِ اعدادِ ۰ تا ۹).

\exercise برنامه‌ای بنویسید که آرایه‌ای از اعداد را \emph{وارونه} چاپ کند:
از خانهٔ آخر تا خانهٔ اول. مراقبِ خطای یکی کم و زیاد باشید: اولین اندیسی
که چاپ می‌کنید چیست؟

\exercise روالی بنویسید که یک آرایه و یک عددِ هدف بگیرد و تعدادِ دفعه‌های
حضورِ هدف در آرایه را برگرداند. با آرایهٔ
\code{[3, 7, 3, 3, 1]} و هدفِ \code{3} بیازمایید (باید ۳ شود).
